\documentclass[conference]{IEEEtran}
\usepackage{amsmath,amssymb,amsthm}
\usepackage{graphicx}
\usepackage{booktabs}
\usepackage{url}
\usepackage{hyperref}
\usepackage{xcolor}

%% --- generated macros ---
\input{../results/macros09}

%% --- theorem environments ---
\newtheorem{lemma}{Lemma}
\newtheorem{theorem}{Theorem}
\newtheorem{corollary}{Corollary}
\newtheorem{definition}{Definition}

% ----------------------------------------------------------------
\begin{document}

\title{NATbreak: Exploiting Enterprise Network Address Translation\\
as a Zero-Cost Adversarial Amplifier\\
for ML-Based Intrusion Detection}

\author{
  \IEEEauthorblockN{Souradeepta Biswas}
  \IEEEauthorblockA{American Express Global Services\\
  New York, NY, USA\\
  sdb.svnit@gmail.com}
}

\maketitle

% ----------------------------------------------------------------
\begin{abstract}
ML-based network intrusion detection systems trained on public datasets
(CICIDS-2017, UNSW-NB15) assume routable per-host IP addresses.
Enterprise networks deploy large-scale NAT: \SifiNEndpoints{} endpoints
sharing \SifiNExternalIps{} external IPs at a \NatRatioLarge{}:1 ratio
in our evaluation environment.
This NAT topology collapses \NIPFeatures{} of \NTotalFeatures{} ML flow
features, degrading Random Forest accuracy from \RfBaselineAcc{} (baseline)
to \TheoremOneRfNatAcc{} under our enterprise NAT configuration.
We introduce NATbreak—three attack variants that exploit NAT topology
awareness (requiring no gradient access to the target model) to further
suppress detection.
We then propose NAT-invariant feature engineering that removes IP-dependent
features and recovers accuracy to \InvRfAccAtSifi{} even at
\NatRatioLarge{}:1 NAT ratio, outperforming the standard RF
by \AccGainInvariant{} accuracy points.
\end{abstract}

\begin{IEEEkeywords}
intrusion detection, adversarial machine learning, NAT, network traffic
classification, feature engineering, CICIDS, UNSW-NB15
\end{IEEEkeywords}

% ----------------------------------------------------------------
\section{Introduction}

Public IDS datasets assign each host a unique IP address.
Enterprise NAT collapses thousands of hosts to a handful of external IPs,
destroying the IP-layer feature space that most classifiers implicitly depend
upon.
This is not a known adversarial technique—it is normal network infrastructure—
but an adversary who knows the target deploys NAT (publicly inferable from port
scan patterns~\cite{RFC3022}) can exploit this collapse as a zero-cost
amplifier.

\textbf{Contributions:}
\begin{enumerate}
  \item Quantification of IP entropy collapse under enterprise NAT
    (Lemma~\ref{lem:collapse}, \S\ref{sec:collapse}).
  \item A formal accuracy-under-NAT model (Theorem~\ref{thm:accuracy},
    \S\ref{sec:model}).
  \item Three NATbreak attack variants requiring only network topology knowledge,
    no model access (\S\ref{sec:attack}).
  \item NAT-invariant feature engineering with formal amplification bound
    (Corollary~\ref{cor:amp}, \S\ref{sec:defense}).
  \item Evaluation against RF, LSTM, and XGBoost on CICIDS-2017 and UNSW-NB15
    (\S\ref{sec:eval}).
\end{enumerate}

% ----------------------------------------------------------------
\section{Background}

\subsection{IDS Datasets and IP Assumptions}

CICIDS-2017~\cite{Sharafaldin2018} and UNSW-NB15~\cite{Moustafa2015} are
the dominant evaluation benchmarks.
Both assign unique routable IPs to each simulated host.
Classifiers trained on these datasets learn IP-based flow attribution:
\texttt{src\_ip}, \texttt{dst\_ip}, and derived features (geo-IP, IP pair
uniqueness) carry significant discriminative weight.

\subsection{Enterprise NAT}

RFC 3022~\cite{RFC3022} describes traditional NAT: internal addresses
(RFC 1918 space) are masqueraded behind one or more external IPs.
At SIFI scale—\SifiNEndpoints{} internal endpoints, \SifiNExternalIps{}
external IPs—the NAT ratio is \NatRatioLarge{}:1.
Every flow originating from any of the \SifiNEndpoints{} hosts appears
to originate from one of \SifiNExternalIps{} IPs.

\subsection{Adversarial ML for Network IDS}

Existing adversarial attacks on network IDS require gradient access~\cite{Carlini2017},
white-box model knowledge~\cite{Goodfellow2014}, or
problem-space constraints~\cite{Biggio2013}.
NATbreak requires none of these: only knowledge that the target uses NAT,
which is publicly inferable.

% ----------------------------------------------------------------
\section{The NAT Feature Collapse Problem}
\label{sec:collapse}

\subsection{Affected Features}

We identify \NIPFeatures{} IP-dependent features (indices 0--2 of our
\NTotalFeatures{}-dimensional flow feature vector):
\texttt{src\_ip\_entropy}, \texttt{dst\_ip\_entropy},
\texttt{ip\_pair\_uniqueness}.
The remaining \NInvariantFeatures{} features are NAT-invariant: packet size
distributions, inter-arrival timing, byte rate, TCP flag entropy, payload
entropy, and port entropy.

\subsection{Classifier Vulnerability}

IP feature weights for the three evaluated classifiers:
RF \RfIpWeight{}, LSTM \LstmIpWeight{}, XGBoost \XgbIpWeight{}.
A classifier with high IP weight is vulnerable to NAT collapse;
the NAT-invariant RF uses IP weight 0.00.

\begin{lemma}[IP Entropy Collapse Under NAT]
\label{lem:collapse}
Let $N$ be the number of internal endpoints and $M$ the number of external
NAT IPs.
The maximum IP source entropy before NAT is $H_{\max}(\text{before}) = \log_2 N$.
After NAT, $H_{\max}(\text{after}) = \log_2 M$.
The collapse magnitude is:
\[
  \Delta H = \log_2 N - \log_2 M = \log_2 \frac{N}{M}.
\]
The collapse fraction is $\Delta H / \log_2 N$.
At $N = \SifiNEndpoints{}$, $M = \SifiNExternalIps{}$:
$\Delta H = \LemmaOneCollapseBits{}$~bits
(collapse fraction \LemmaOneCollapseFrac{}).
\end{lemma}

\begin{proof}
The source IP distribution has maximum entropy $\log_2 N$ when all $N$
endpoints are equally likely.
After NAT maps $N$ IPs to $M$ external IPs, the observable distribution
has support size $M$, bounding entropy at $\log_2 M$.
The collapse $\Delta H = \log_2(N/M)$ follows directly.
\end{proof}

Figure~\ref{fig:entropy} illustrates entropy collapse as a function of NAT
ratio.

\begin{figure}[!t]
  \centering
  \includegraphics[width=\columnwidth]{../src/figures/fig3_entropy_collapse.pdf}
  \caption{IP entropy before and after NAT (left) and collapse fraction (right)
    as NAT ratio increases.}
  \label{fig:entropy}
\end{figure}

% ----------------------------------------------------------------
\section{Accuracy-Under-NAT Model}
\label{sec:model}

\begin{theorem}[Classifier Accuracy Under NAT]
\label{thm:accuracy}
Let $A_0$ be baseline accuracy (no NAT), $w_{ip}$ the IP feature weight
fraction, and $r = N/M$ the NAT ratio.
The expected accuracy under NAT is:
\[
  A(r) = A_0 \cdot \left(1 - w_{ip} \cdot \left(1 - \frac{1}{r}\right)\right).
\]
$A(r)$ is monotone non-increasing in $r$.
As $r \to \infty$: $A \to A_0(1 - w_{ip})$.
At $r = 1$ (no NAT): $A = A_0$.
For RF at $r = \NatRatioLarge{}$: $A = \TheoremOneRfNatAcc{}$.
\end{theorem}

\begin{proof}
The IP features contribute fraction $w_{ip}$ of total discriminative power.
Under NAT ratio $r$, IP entropy collapses by factor $(1 - 1/r)$ (Lemma~\ref{lem:collapse}),
reducing their contribution proportionally.
The NAT-invariant features ($1 - w_{ip}$ fraction) are unaffected.
Multiplying through gives $A(r)$ as stated.
Monotonicity follows because $(1 - 1/r)$ is monotone increasing in $r$.
\end{proof}

Figure~\ref{fig:nat_acc} shows the model prediction and the NAT-invariant
classifier's constant accuracy across NAT ratios.

\begin{figure}[!t]
  \centering
  \includegraphics[width=\columnwidth]{../src/figures/fig2_accuracy_vs_nat_ratio.pdf}
  \caption{RF accuracy vs NAT ratio: baseline, NAT-degraded, and NAT-invariant.}
  \label{fig:nat_acc}
\end{figure}

% ----------------------------------------------------------------
\section{NATbreak Attack}
\label{sec:attack}

\textbf{Adversary model:} The attacker knows the target network uses NAT
(inferred from scan patterns) but has no access to the IDS model, its weights,
or training data.
Attack flows originate from inside the NAT perimeter.

\subsection{Type I: IP Feature Pollution}

Generate a high volume of benign-looking cover flows from the same external
NAT IP as the attack flows.
With pollution ratio \TypeOnePollutionRatio{}:1 (cover:attack), the IP-based
score of the attack flows is diluted toward the benign cluster.
The IP feature contribution to the malicious score is reduced by:
$1/(1 + \text{pollution\_ratio})$.

\subsection{Type II: Flow Tuple Collision}

Time attack flows to share the NAT state table entry of a legitimate
concurrent flow.
The IDS observes the attack flow's packet as belonging to the legitimate
flow's 5-tuple, inheriting its benign IP features.
Requires a \TypeTwoCollisionWindowMs{}~ms timing window.

\subsection{Type III: Feature Space Saturation}

Push IP feature values toward the classifier's decision boundary (0.5) using
knowledge of the NAT ratio only.
Saturation fraction: \TypeThreeSaturationFrac{}.
This makes the attack flow's IP-based score ambiguous without affecting
NAT-invariant features.

Figure~\ref{fig:attacks} shows accuracy under each attack variant for all
three classifiers.

\begin{figure}[!t]
  \centering
  \includegraphics[width=\columnwidth]{../src/figures/fig4_attack_by_variant.pdf}
  \caption{Classifier accuracy under each NATbreak variant at \NatRatioLarge{}:1 NAT.}
  \label{fig:attacks}
\end{figure}

% ----------------------------------------------------------------
\section{NAT-Invariant Defense}
\label{sec:defense}

\subsection{Feature Engineering}

Remove the \NIPFeatures{} IP-dependent features from the flow representation.
The reduced \NInvariantFeatures{}-dimensional vector uses only:
flow duration, packet size statistics, inter-arrival timing, byte rate,
TCP flag entropy, payload entropy, and port entropy.

\begin{corollary}[NATbreak Amplification Factor]
\label{cor:amp}
The adversarial amplification factor of NATbreak is $F = N/M$ (the NAT ratio).
For a classifier with IP weight $w_{ip}$, the maximum additional accuracy
degradation achievable by NATbreak beyond bare NAT is bounded by
$w_{ip} \cdot (1 - 1/r)$, which scales with $F$.
For the NAT-invariant classifier ($w_{ip} = 0$), $F = 0$:
NATbreak provides zero amplification regardless of NAT ratio.
At SIFI NAT ratio \NatRatioLarge{}: $F = \CorollaryOneAmpSifi{}$.
\end{corollary}

Figure~\ref{fig:invariant} compares standard RF and NAT-invariant RF under
increasing NAT ratio and NATbreak attack.

\begin{figure}[!t]
  \centering
  \includegraphics[width=\columnwidth]{../src/figures/fig5_invariant_comparison.pdf}
  \caption{Standard vs NAT-invariant RF: accuracy under NAT and NATbreak.}
  \label{fig:invariant}
\end{figure}

% ----------------------------------------------------------------
\section{Evaluation}
\label{sec:eval}

\subsection{Setup}

Classifiers: Random Forest (\RfBaselineAcc{} baseline), LSTM
(\LstmBaselineAcc{}), XGBoost (\XgbBaselineAcc{}), RF-NAT-Invariant
(\InvBaselineAcc{}).
Datasets: CICIDS-2017 (\CicidsNBenign{} benign / \CicidsNMalicious{} malicious
flows) and UNSW-NB15 (\UnswNBenign{} / \UnswNMalicious{}).
NAT configs: no NAT ($r=1$), small ($r=\NatRatioSmall{}$), medium
($r=\NatRatioMedium{}$), large SIFI ($r=\NatRatioLarge{}$).

\subsection{Baseline vs NAT Accuracy}

Figure~\ref{fig:baseline} shows the accuracy gap between baseline and NAT
conditions for all four classifiers at the SIFI NAT ratio.
RF degrades from \RfBaselineAcc{} to \NatAccRF{};
LSTM from \LstmBaselineAcc{} to \NatAccLSTM{};
XGBoost from \XgbBaselineAcc{} to \NatAccXGBoost{}.
The NAT-invariant RF maintains \NatAccRFNATInvariant{}, demonstrating
immunity to NAT collapse.

\begin{figure}[!t]
  \centering
  \includegraphics[width=\columnwidth]{../src/figures/fig1_baseline_vs_nat.pdf}
  \caption{Classifier accuracy: no-NAT baseline vs SIFI NAT ($r=\NatRatioLarge{}$).}
  \label{fig:baseline}
\end{figure}

\subsection{NATbreak Attack Results}

At \NatRatioLarge{}:1 NAT ratio, NATbreak Type~I further reduces RF accuracy
to \RfAttackTypeOneAcc{}; Type~II to \RfAttackTypeTwoAcc{}; Type~III to
\RfAttackTypeThreeAcc{}.
The NAT-invariant RF is unaffected by all three variants.

\subsection{Amplification Factor}

Figure~\ref{fig:amp} shows the NATbreak amplification factor scaling
linearly with NAT ratio, reaching $F = \AmpFactorLarge{}$ at
SIFI scale vs $F = \AmpFactorSmall{}$ for small enterprise.

\begin{figure}[!t]
  \centering
  \includegraphics[width=\columnwidth]{../src/figures/fig6_amplification.pdf}
  \caption{NATbreak amplification factor $F = N/M$ vs NAT ratio (log-log).}
  \label{fig:amp}
\end{figure}

% ----------------------------------------------------------------
\section{Related Work}

Table~\ref{tab:related} compares NATbreak to existing adversarial IDS attacks.

\begin{table}[!t]
  \centering
  \caption{Adversarial IDS Attack Comparison}
  \label{tab:related}
  \begin{tabular}{lccc}
    \toprule
    Attack & Model Access & NAT Aware & Zero Cost \\
    \midrule
    FGSM~\cite{Goodfellow2014} & Gradient & -- & -- \\
    C\&W~\cite{Carlini2017} & White-box & -- & -- \\
    Biggio~\cite{Biggio2013} & White-box & -- & -- \\
    NATbreak (ours) & None & \checkmark & \checkmark \\
    \bottomrule
  \end{tabular}
\end{table}

\textbf{Dataset bias:} Sommer and Paxson~\cite{Sommer2010} identified the
dataset-deployment gap as a fundamental IDS challenge.
Our work identifies NAT as a specific, quantifiable, and exploitable
instance of this gap.
\textbf{Evasion attacks:} Rigaki and Garcia~\cite{Rigaki2018} and
Rosenberg et al.~\cite{Rosenberg2021} demonstrate evasion with GAN-generated
and query-based attacks respectively; both require model interaction.
\textbf{Feature selection:} Buczak and Guven~\cite{Buczak2016} survey
feature engineering for IDS without addressing NAT invariance.

% ----------------------------------------------------------------
\section{Limitations and Future Work}

Theorem~\ref{thm:accuracy} models accuracy as linear in the IP collapse term.
Real classifiers may exhibit non-linear interactions between IP and invariant
features; calibrating the model against Snort~\cite{Snort} and
Suricata~\cite{Suricata} rule engines is future work.
Type~II (tuple collision) depends on timing precision; success probability
as a function of connection rate is not modeled here.
Extending NATbreak to IPv6 environments with NAT64 and CLAT translation
layers is a natural extension.

% ----------------------------------------------------------------
\section{Conclusion}

Enterprise NAT is a ubiquitous structural property that collapses
\LemmaOneCollapseBits{}~bits of IP entropy at SIFI scale and degrades
standard RF IDS accuracy from \RfBaselineAcc{} to \TheoremOneRfNatAcc{}.
NATbreak exploits this collapse as a zero-cost adversarial amplifier—
no model access required—further suppressing detection.
NAT-invariant feature engineering eliminates the IP dependency, recovering
\InvRfAccAtSifi{} accuracy at \NatRatioLarge{}:1 NAT ratio and rendering
NATbreak ineffective.

% ----------------------------------------------------------------
\bibliographystyle{IEEEtran}
\bibliography{paper09}

\end{document}
